\section{Introduction}
\label{sec:introduction}
\PARstart{T}{he} Industrial Internet of Things (IIoT) connects operational technology devices, sensors, programmable logic controllers, industrial cameras, and edge gateways to IP-based networks, supporting automation, process visibility, and remote monitoring in sectors such as manufacturing, energy, healthcare, and logistics. This increased connectivity also requires industrial networks to process traffic generated by heterogeneous devices and to identify anomalous behavior without disrupting operational continuity \cite{firouzi2025datasense,esmaeilyfard2026gelax}.

Network intrusion detection systems (IDSs) in IIoT environments must analyze traffic produced by devices with different communication patterns, sampling rates, protocols, and operational roles. Although supervised classifiers can provide effective detection when representative labels are available, obtaining and maintaining such labels may be difficult in evolving industrial deployments. Device populations, network configurations, and traffic behavior may change over time, whereas reliable annotation generally requires domain expertise and controlled investigation. These conditions motivate methods that can organize and refine IIoT traffic features before comprehensive attack labels are available \cite{andresini2021insomnia,alrumaih2026lightweight}.

Unsupervised feature selection (UFS) addresses this setting by identifying informative attributes directly from unlabeled data. For IIoT data, this process should consider that discriminative behavior may arise from interactions among packet sizes, header flags, timing measurements, address diversity, port usage, fragmentation indicators, and traffic-rate statistics, rather than from the independent relevance of individual features \cite{solorio2020review}.

Accordingly, feature selection in this work is not restricted to aggressive dimensionality reduction. We distinguish between two operating regimes: compact selection, in which a strict feature budget is imposed, and representation-aware feature refinement, in which the objective is to preserve the behavioral structure of the original traffic while removing attributes that do not improve that structure. \rev{This distinction is stated here once and is used throughout the paper.}

Classical UFS methods, such as Laplacian Score \cite{he2005laplacian}, SPEC \cite{zhao2007spectral}, MCFS \cite{cai2010unsupervised}, UDFS \cite{yang2011l2}, and NDFS \cite{li2012unsupervised}, provide efficient mechanisms for ranking or selecting features without labels. These methods generally preserve local, spectral, or graph-based structures in the original feature space. However, when relevant traffic behavior depends on nonlinear interactions among correlated attributes, independently ranking original features may not fully represent the structure needed for subsequent clustering or classification. This motivates evaluating candidate subsets according to the behavioral representation they preserve, rather than only according to their geometry in the raw input space.

Self-supervised representation learning provides a mechanism for learning such a behavioral space from unlabeled traffic. SimCLR learns representations by contrasting positive and negative examples \cite{chen2020simple}, whereas VICReg is a non-contrastive joint-embedding method that avoids negative pairs and combines invariance, variance, and covariance regularization \cite{bardes2022vicreg}. In tabular IIoT data, however, the construction of augmented views requires particular care: arbitrary feature corruption may disrupt combinations of attributes that jointly describe a network behavior. We therefore organize the traffic features into semantic Context Groups and use these groups to define domain-aware augmentations. The groups are based on the measurement semantics of the attributes, rather than on attack labels, and represent categories such as packet size, header flags, timing, traffic rate, address diversity, and fragmentation.

In this paper, we propose \method{}, an unsupervised pipeline for representation-aware feature refinement and cluster interpretation in IIoT intrusion detection. As illustrated in Fig.~\ref{fig:overview}, the method combines self-supervised representation learning, latent-space subset evaluation, and post-hoc explainability. Class labels are not used to train the encoder, select features, or construct the latent clusters; they are introduced only during downstream evaluation. The main contributions are summarized as follows:

\textbf{C1 --- Group-Aware VICReg Representation Learning.} We train a VICReg encoder using a group-aware augmentation strategy in which semantic Context Group masking is the primary domain-aware operation, complemented by low-amplitude Gaussian noise and individual feature dropout \cite{bardes2022vicreg}.

\textbf{C2 --- Latent-Space Bidirectional Feature Refinement.} We introduce a bidirectional Silhouette-based search that evaluates original-feature subsets through the representation produced by a frozen VICReg encoder. The search is initialized with the complete feature set and alternates backward removal with forward recovery of previously removed attributes. The Silhouette score is used as a label-free surrogate for latent cluster cohesion and separation; it does not guarantee optimal downstream classification.

\textbf{C3 --- SHAP-Based Cluster Profiling.} We use a Random Forest proxy to approximate the cluster assignments obtained in the latent space and apply TreeExplainer SHAP to characterize the contribution of the selected original traffic features to each proxy-predicted cluster \cite{lundberg2017unified}. \rev{The fidelity of the proxy is measured on held-out data, and a quantitative clarity index determines which cluster profiles are concentrated and faithful enough to be interpreted.}

\Figure[t!](topskip=0pt, botskip=0pt, midskip=0pt)[width=0.95\textwidth]{figures/fig1_overview.png}
{Overview of the proposed \method{} pipeline and its main contributions.\label{fig:overview}}

Experiments are conducted on the DataSense CIC IIoT 2025 dataset using 10-fold stratified cross-validation \cite{firouzi2025datasense}. Within each fold, the scaler is fitted on the training partition, the VICReg encoder is trained using only that partition, feature selection is performed on the corresponding training representation, and the downstream classifier is evaluated on the untouched test partition. \rev{At its natural operating point, the method retains on average 59 of the 71 numeric traffic features and achieves macro-F1 scores of 0.936 for binary detection and 0.908 for eight-class attack-category classification with Random Forest. A paired equivalence test confirms that this performance lies within $\pm$0.01 macro-F1 of the all-feature configuration, and the method significantly outperforms MCFS and the supervised DataSense-17 subset. The evaluation is complemented by three independent master seeds, device- and attack-execution-grouped folds, controlled component comparisons, and two additional public IoT/IIoT datasets, N-BaIoT and WUSTL-IIoT-2021.}

Overall, \method{} is intended for settings in which preserving unlabeled behavioral structure and supporting subsequent interpretation are more important than minimizing feature cardinality alone. Its cluster profiles\rev{, which are stable across folds and seeds and are reproduced by the proxy on unseen data,} provide behavioral evidence related to traffic characteristics such as TCP-flag activity, time-to-live statistics, and packet-size distributions, while remaining independent of attack labels during representation learning and feature selection.


The remainder of this paper is organized as follows. Section~\ref{sec:related} reviews related work. Section~\ref{sec:dataset} describes the DataSense CIC IIoT 2025 dataset and its feature organization. Section~\ref{sec:method} presents the \method{} pipeline. Section~\ref{sec:setup} details the experimental protocol\rev{, including the grouped protocols, the independent master seeds, and the external dataset}. Section~\ref{sec:results} reports the results. Section~\ref{sec:discussion} discusses implications and limitations, and Section~\ref{sec:conclusion} concludes the paper.
